% !TEX root = ../../cookbook.tex
\newpage
\recipe{Birria de Calzonso}{Para 6--8 porciones. Puedes consultar el \href{https://www.youtube.com/watch?v=g2F5RO6vNSs}{video de la receta de ``La Capital'' en YouTube}.}{}{}

\vspace*{-0.6cm}
\subsection*{Ingredientes}
\vspace*{-0.4cm}
\begin{multicols}{2}

    \begin{ingredients}
        \item 2 kg de carne de res\footnote{Usar carne con partes magras y grasas: chamorro, diezmillo, costillas, rabo de res.}
        \item 2--3 huesos de tuétano
        \item 7 chiles guajillo\footnote{\label{note:chiles}Secos, desvenados, sin semillas.}
        \item 5 chiles ancho\footref{note:chiles}
        \item Chiles de árbol, al gusto\footnote{Estos chiles aportarán el picor.}
        \item 1 cebolla mediana, en cuartos
        \item 6 dientes de ajo
        \item 1 pedazo pequeño de jengibre, de aproximadamente 2 cm
        \item 1 cucharadita de comino molido
        \item 1 cucharadita de orégano seco
        \item \sfrac{1}{2} cucharadita de clavo molido
        \item \sfrac{1}{2} cucharadita de pimienta negra
        \item 3 hojas de laurel
        \item 1 vara de canela
        \item Amor (imprescindible)
        \item Chiles morita
        \item Queso Oaxaca\footnote{Si no encuentra queso Oaxaca, se puede mozzarella.}
        \item Sal
        \item Aceite de oliva
        \item Tortillas de maíz
        \item Cilantro y cebolla frescos
    \end{ingredients}
\end{multicols}

\vspace{-0.5cm}

% --- Preparación ---

Colocar los \textbf{chiles guajillo}, \textbf{chiles ancho} y \textbf{chiles de árbol} en una olla y cúbralos con agua, prenda el fuego y cuando empiece a hervir, apague el fuego y deje reposar aproximadamente 20 minutos. 
Escurra los chiles.
En un sartén seco, a fuego bajo, tostar el \textbf{ajo}, \textbf{jengibre}, \textbf{comino}, \textbf{orégano}, \textbf{clavo}, \textbf{pimienta negra} y \textbf{canela}. 
Procure tostar cada especia por separado, ya que necesitan diferentes tiempos de cocción, procura no quemar las especias porque pueden amargar.
Licue los \textbf{chiles} ablandados, \textbf{cebolla}, junto con las especies que tostamos hasta obtener una mezcla homogénea.
Coloque la \textbf{carne de res} en una olla grande y cúbrala con el \textbf{adobo} (la mezcla licuada). Marine la carne por al menos 2 horas.

Agregue \textbf{agua} hasta cubrir la carne, junto con los \textbf{huesos de tuétano} y las \textbf{hojas de laurel}. 
Cocine a fuego bajo durante 3--4 horas, o hasta que la carne esté suave. Añada la \textbf{sal} únicamente después de la tercera hora.
Lave y seque los \textbf{chiles morita}. 
Caliente \textbf{aceite de oliva} en un sartén pequeño. Fría los chiles hasta que se inflen, pero sin que se quemen. 
Déjelos enfriar y añádalos al caldo junto con el aceite.
Una vez que la birria esté lista, retire la capa de grasa de la superficie del caldo. Reserve la grasa para cocinar las quesadillas.

Para preparar las quesadillas, sumerja una \textbf{tortilla de maíz} en la grasa reservada y colóquela en un sartén caliente. 
Añada la \textbf{carne de res} y el \textbf{queso}, doble la tortilla y cocínela hasta que el queso se derrita y la tortilla esté ligeramente crujiente.
Sirva la birria con las quesadillas, \textbf{cilantro} y \textbf{cebolla} picados. Sirva el caldo por separado en una taza.

% --- Descripción ---

\medskip